\chapter{Anexos}
\label{chap:anexos}

\section{Fragmentos de código}

\begin{lstlisting}[float=htbp, style=CSharp, caption={Implementación del patrón Static Factory Method en la entidad User.}, label={lst:factory_user}]
public sealed class User : Entity
{
    // Constructor privado: impide la instanciacion directa y anula estados invalidos
    private User(Guid id, string firstName, string lastName, string email) { ... }

    // Factory Method estatico: Valida invariantes y emite el evento de dominio
    public static User Create(string firstName, string lastName, string email)
    {
        User user = new(Guid.NewGuid(), firstName, lastName, email);
        user.Raise(new UserCreatedDomainEvent(
            user.Id, user.FirstName, user.LastName, user.Email));

        return user;
    }
}
\end{lstlisting}

\begin{lstlisting}[float=htbp, style=CSharp, caption={Implementación del patrón Decorator mediante Pipeline Behaviors para trazabilidad y validación.}, label={lst:decorator_mediatr}]
// Logging: mide y traza todas las peticiones
public sealed class LoggingBehavior<TRequest, TResponse>(ILogger<...> logger)
    : IPipelineBehavior<TRequest, TResponse>
{
    public async Task<TResponse> Handle(TRequest request, RequestHandlerDelegate<TResponse> next, CancellationToken ct)
    {
        logger.LogInformation("[{Behavior}] Executing: {Name}", 
            nameof(LoggingBehavior<,>), typeof(TRequest).Name);
            
        Stopwatch timer = Stopwatch.StartNew();
        TResponse response = await next(ct);
        
        logger.LogInformation("[{Behavior}] Completed in {Ms}ms", 
            nameof(LoggingBehavior<,>), timer.ElapsedMilliseconds);
        return response;
    }
}

// Validacion: ejecuta todos los IValidator<TRequest> registrados
public sealed class ValidationBehavior<TRequest, TResponse>(
    IEnumerable<IValidator<TRequest>> validators)
    : IPipelineBehavior<TRequest, TResponse>
{
    public async Task<TResponse> Handle(TRequest request, RequestHandlerDelegate<TResponse> next, CancellationToken ct)
    {
        if (!validators.Any()) return await next(ct);
        
        ValidationResult[] results = await Task.WhenAll(
            validators.Select(v => v.ValidateAsync(new(request), ct)));
            
        List<ValidationFailure> errors = results
            .Where(r => !r.IsValid)
            .SelectMany(r => r.Errors).ToList();
            
        if (errors.Count > 0) throw new ValidationException(errors);
        return await next(ct);
    }
}
\end{lstlisting}

\begin{lstlisting}[float=htbp, style=TypeScript, caption={Implementación del patrón Proxy mediante interceptores HTTP para inyección y rotación de tokens.}, label={lst:proxy_interceptors}]
// Request: inyecta el token de acceso en cada peticion
apiClient.interceptors.request.use((config) => {
  const requestConfig = config as RetriableRequestConfig;
  if (!requestConfig.skipAuthorization) {
    const accessToken = getAccessToken();
    if (accessToken) {
        requestConfig.headers.set('Authorization', `Bearer ${accessToken}`);
    }
  }
  return requestConfig;
});

// Response: detecta 401, refresca y reintenta (una sola vez)
apiClient.interceptors.response.use(
  (response) => response,
  async (error: AxiosError) => {
    const requestConfig = error.config as RetriableRequestConfig | undefined;
    if (!requestConfig || requestConfig._retry || error.response?.status !== 401) {
        return Promise.reject(error);
    }
    
    requestConfig._retry = true;
    // Promesa compartida para sincronizar multiples peticiones fallidas
    const refreshedToken = await getRefreshTokenOnce();  
    requestConfig.headers.set('Authorization', `Bearer ${refreshedToken}`);
    
    return apiClient(requestConfig);
  }
);
\end{lstlisting}

\begin{lstlisting}[float=htbp, style=TypeScript, caption={Implementación del patrón Facade para agrupar los servicios de llamadas HTTP.}, label={lst:facade_apigateway}]
export const ApiGateway = {
  organization: {
    invitations: OrganizationInvitationsGateway,
    management:  OrganizationManagementGateway,
    members:     OrganizationMembersGateway,
    roles:       OrganizationRolesGateway,
  },
  project: {
    customFields: ProjectCustomFieldsGateway,
    management:   ProjectManagementGateway,
  },
  workItem: {
    management:  WorkItemsGateway,
    workLogs:    WorkLogGateway,
    attachments: AttachmentGateway,
    comments:    CommentGateway,
  },
  user:  { session: SessionGateway, users: UsersGateway },
  agent: AgentGateway,
  chat:  ChatGateway,
} as const;

// Uso simplificado desde cualquier componente:
// ApiGateway.workItem.management.createWorkItem(projectId, payload)
\end{lstlisting}

\begin{lstlisting}[float=htbp, style=CSharp, caption={Definición de interfaces base para CQRS y su aplicación práctica.}, label={lst:cqrs_pattern}]
public interface ICommand : IRequest;
public interface ICommand<out TResponse> : IRequest<TResponse>;
public interface IQuery<out TResponse> : IRequest<TResponse>;

// Command: modifica el estado del dominio
public sealed record CreateOrganizationCommand(
    string Name, string Description, string Website)
    : ICommand<Result<OrganizationDto>>;

// Query: operacion de solo lectura
public sealed record GetMyOrganizationsQuery(
    int Limit, string? Cursor, IReadOnlyCollection<string> Sort)
    : IQuery<Result<PaginatedResult<SimpleOrganizationDto>>>;
\end{lstlisting}

\begin{lstlisting}[float=htbp, style=CSharp, caption={Uso del patrón Mediator para desacoplar controladores de sus manejadores.}, label={lst:mediator_pattern}]
// Controller base que inyecta el ISender del Mediator
public abstract class ApiBaseController : ControllerBase
{
    private ISender? _sender;
    protected ISender Sender => _sender ??= 
        HttpContext.RequestServices.GetRequiredService<ISender>();
}

// Controller concreto enviando un Command
[HttpPost]
public async Task<IActionResult> Create(
    [FromBody] CreateOrganizationRequest request, CancellationToken ct)
{
    var command = new CreateOrganizationCommand(
        request.Name, request.Description, request.Website);
    return HandleResult(await Sender.Send(command, ct));
}
\end{lstlisting}

\begin{lstlisting}[float=htbp, style=CSharp, caption={Separación de contrato e implementación mediante el patrón Repository.}, label={lst:repository_pattern}]
// Dominio: contrato agnostico a la infraestructura
public interface IOrganizationRepository
{
    Task<Organization?> GetByIdAsync(Guid id, CancellationToken ct = default);
    Task AddAsync(Organization org, CancellationToken ct = default);
    Task SaveChangesAsync(CancellationToken ct = default);
}

// Infraestructura: implementacion concreta acoplada a EF Core
public class OrganizationRepository(OrganizationDbContext dbContext) 
    : IOrganizationRepository
{
    public async Task<Organization?> GetByIdAsync(Guid id, CancellationToken ct = default) =>
        await dbContext.Organizations
            .AsSplitQuery()
            .Include(o => o.Invitations)
            .Include(o => o.Roles).ThenInclude(r => r.Permissions)
            .FirstOrDefaultAsync(o => o.Id == id, ct);

    public async Task SaveChangesAsync(CancellationToken ct = default) =>
        await dbContext.SaveChangesAsync(ct);
}
\end{lstlisting}

\begin{lstlisting}[float=htbp, style=CSharp, caption={Implementación del patrón Observer para Eventos de Dominio.}, label={lst:observer_domain}]
public abstract class Entity
{
    private readonly List<IDomainEvent> _domainEvents = [];
    public IReadOnlyCollection<IDomainEvent> DomainEvents => _domainEvents.AsReadOnly();
    public void Raise(IDomainEvent domainEvent) => _domainEvents.Add(domainEvent);
}

// Intercepcion del DbContext para despachar observadores
public override async Task<int> SaveChangesAsync(CancellationToken ct = default)
{
    List<IDomainEvent> domainEvents = ChangeTracker.Entries<Entity>()
        .SelectMany(e => e.Entity.DomainEvents).ToList();

    await domainEventsDispatcher.DispatchAsync(domainEvents, ct);
    InsertOutboxMessages(domainEvents);
    
    return await base.SaveChangesAsync(ct);
}
\end{lstlisting}

\begin{lstlisting}[float=htbp, style=CSharp, caption={Consumo asíncrono e idempotente de Eventos de Integración (Inbox).}, label={lst:observer_integration}]
public abstract class IntegrationEventHandler<TIntegrationEvent, TContext>(...) 
    : INotificationHandler<TIntegrationEvent>
{
    public async Task Handle(TIntegrationEvent notification, CancellationToken ct)
    {
        Guid messageId = integrationEventContextAccessor.CurrentMessageId 
            ?? GetFallbackMessageId(notification);

        // Verificacion de idempotencia: actua como un observer seguro
        bool alreadyProcessed = await dbContext.Set<InboxMessage>()
            .AsNoTracking().AnyAsync(m => m.Id == messageId, ct);
            
        if (alreadyProcessed) return;

        await HandleIntegrationEvent(notification, ct);

        dbContext.Set<InboxMessage>().Add(InboxMessage.Create(messageId, ...));
        await dbContext.SaveChangesAsync(ct);
    }
}
\end{lstlisting}

\begin{lstlisting}[float=htbp, style=TypeScript, caption={Suscripción reactiva al estado del servidor mediante React Query.}, label={lst:observer_frontend}]
// Estructura jerarquica de claves para invalidacion granular
export const QueryKeys = {
  workItem: {
    management: {
      all: ['workitem', 'management'] as const,
      list: (projectId: string) => ['workitem', 'management', 'list', projectId] as const,
    },
  },
} as const;

// Hook observador del estado
export const useProjectWorkItemsQuery = (projectId: string | null) => {
  const accessToken = useAuthStore((state) => state.accessToken);
  return useQuery({
    queryKey: QueryKeys.workItem.management.list(projectId ?? ''),
    queryFn: async () => { /* Logica de paginacion y llamada a la API */ },
    enabled: Boolean(accessToken && projectId),
  });
};
\end{lstlisting}

\begin{lstlisting}[float=htbp, style=TypeScript, caption={El patrón Command encapsulando operaciones de escritura en React Query.}, label={lst:command_frontend}]
export const useCreateWorkItemMutation = (projectId: string) => {
  const handleApiError = useApiErrorHandler();

  return useMutation({
    mutationKey: [...QueryKeys.workItem.management.all, 'create', projectId],
    mutationFn: (payload: CreateWorkItemRequest) =>
      ApiGateway.workItem.management.createWorkItem(projectId, payload),
    onSuccess: async () => {
      // Efectos secundarios unificados dentro del Command
      await queryClient.invalidateQueries({ 
        queryKey: QueryKeys.workItem.management.list(projectId) 
      });
      toast.success(i18n.t('create.feedback.created', { ns: 'workitems' }));
    },
    onError: (error) => handleApiError(error),
  });
};
\end{lstlisting}

\begin{lstlisting}[float=htbp, style=CSharp, caption={Implementación de Strategy para la inyección de datos iniciales.}, label={lst:strategy_pattern}]
public interface IUserSeederStrategy
{
    Task SeedAsync(UserDbContext dbContext, CancellationToken ct = default);
}

public sealed class DevelopmentUserSeeder : IUserSeederStrategy
{
    public async Task SeedAsync(UserDbContext dbContext, CancellationToken ct = default)
    {
        if (await dbContext.Users.AnyAsync(ct)) return;
        User[] users = [
            User.Create("Admin", "Developer", "admin@dev.com"),
            User.Create("John",  "Doe",       "john@dev.com"),
        ];
        await dbContext.Users.AddRangeAsync(users, ct);
        await dbContext.SaveChangesAsync(ct);
    }
}

// Inyeccion de dependencias resolviendo la estrategia en el arranque
services.AddScoped<IUserSeederStrategy>(
    isDevelopment 
        ? _ => new DevelopmentUserSeeder() 
        : _ => new ProductionUserSeeder());
\end{lstlisting}

\begin{lstlisting}[float=htbp, style=CSharp, caption={Implementación del patrón Factory delegando logicas de expiracion.}, label={lst:factory_token}]
public class RefreshTokenFactory(IRefreshTokenExpirationPolicy policy) : IRefreshTokenFactory
{
    public RefreshToken CreateForUser(Guid userId)
    {
        // La fabrica obtiene la politica de expiracion inyectada dinamicamente
        TimeSpan duration = policy.ExpirationDuration;
        return RefreshToken.Create(userId, duration);
    }
}
\end{lstlisting}

\begin{lstlisting}[float=htbp, style=CSharp, caption={Sobrescritura de SaveChangesAsync para persistir eventos de dominio.}, label={lst:savechanges}]
public override async Task<int> SaveChangesAsync()
{
    List<IDomainEvent> domainEvents = GetDomainEvents();

    await domainEventsDispatcher.DispatchAsync(domainEvents, cancellationToken);

    InsertOutboxMessages(domainEvents);

    return await base.SaveChangesAsync(cancellationToken);
}
\end{lstlisting}

\begin{lstlisting}[float=htbp, style=CSharp, caption={Implementación del patrón Specification para la capa de persistencia de datos.}, label={lst:specification_pattern}]
public abstract class BaseSpecification<T>(Expression<Func<T, bool>> criteria) 
    : ISpecification<T>
{
    private readonly List<Expression<Func<T, object>>> _includes = [];
    public Expression<Func<T, bool>> Criteria { get; } = criteria;
    public IReadOnlyCollection<Expression<Func<T, object>>> Includes 
        => _includes.AsReadOnly();
        
    protected void AddInclude(Expression<Func<T, object>> expr) 
        => _includes.Add(expr);
}

// Especificacion concreta que encapsula la logica de filtrado e inclusion
public sealed class UserWithActiveRefreshTokensSpecification 
    : BaseSpecification<User>
{
    public UserWithActiveRefreshTokensSpecification(Guid userId)
        : base(user => user.Id == userId)
    {
        AddInclude(user => user.RefreshTokens
            .Where(rt => rt.RevokedAt == null && rt.ExpiresAt > DateTime.UtcNow));
    }
}
\end{lstlisting}

\begin{lstlisting}[float=htbp, style=CSharp, caption={Validación en cascada mediante Chain of Responsibility.}, label={lst:permission_validation}]
public sealed class OrganizationPermissionBehavior<TRequest, TResponse>(...) 
    : IPipelineBehavior<TRequest, TResponse>
{
    public async Task<TResponse> Handle(TRequest request, 
        RequestHandlerDelegate<TResponse> next, CancellationToken ct)
    {
        // [...] Extraccion de la configuracion de permisos requeridos
        
        Organization? org = await organizationRepository
            .GetByIdAsync(permissionRequest.OrganizationId, ct);
        if (org is null) return CreateFailureResponse(OrganizationNotFound); // Corta la cadena

        OrganizationInvitation? invitation = org.Invitations
            .FirstOrDefault(inv => inv.UserId == userContext.UserId);
        if (invitation is null) return CreateFailureResponse(OrganizationMembershipRequired);

        // [...] Evaluacion del nivel del rol
        if (rolePermission?.Level == PermissionLevel.Full) 
            return await next(ct); // Autorizado: pasa el control al siguiente eslabon
            
        return CreateFailureResponse(OrganizationPermissionDenied);
    }
}
\end{lstlisting}